\section{概述}
\begin{enumerate}
  \item 以理解为主。反复背诵。反复做题。
  \item \zx{应试时候，不要添加题干没有的信息，题目没有明确强调，则应当 }。
  \item 行政法核心是\zd{控权}。控权方式，\zd{合法行政}。\bs{合法的法}，背。
        \begin{enumerate}
          \item 法律。全人大/全人常。xx 法。
          \item 法规 = 行政法规（\zd{国务院}）+ 地方性法规（\zd{省、市人大人常}）。\zd{条例}、办法、规定。
          \item 规章 = 部门规章（\zd{组基直专}）+ 地方性规章（\zd{省、市政府}）。办法、规定。\bs{不能叫条例}。
          \item 其他规范性文件。俗称\zd{红头文件}，谁都可以制定。通知、意见、公告。
          \item 命令序列号。\bs{规章以上带令字编号}。\al{农业部二零一五年第五号令}。
        \end{enumerate}
  \item 官管民，民告官。官 = 行政机关、授权组织。管 = 管理行为。民 = 相对人、相关人。告 = 诉讼、复议、国家赔偿。官 = 被告。
        \begin{kk}{6 个基本原则。必考，背}
          \begin{enumerate}
            \item 合法行政。是\zd{形式法}（成文法）的体现。\bs{合法。优先、保留；抵触、授权、先后}。
                  \begin{enumerate}
                    \item 法律优先。\bs{不得抵触}。\al{规定只能罚 5000，罚了 8000}。
                    \item 法律保留。\bs{无授权无行政，先立法后执法}。
                    \item vs 依法行政（上位概念）。\al{我的手是我的一个补充 \cuo。我的手是我的一个部分}。
                    \item \zx{vs 合理行政}。\al{合理行政是合法行政的必然要求 \cuo{延伸}}。
                  \end{enumerate}
            \item 合理行政。是\zx{实质法}（自然法，精神、理念）的体现。\zd{合乎常理}。\bs{合理。公、相、比}。
                  \begin{enumerate}
                    \item 公平公正（平等）。\zd{同等情况同等对待；差别情况差别对待}（实质平等）。
                    \item 相关原则。要有正当的、实质的\zd{关联}。
                    \item 比例原则。\zd{少侵害}。\bs{合、适、必；益、手、侵}。
                          \begin{enumerate}
                            \item 合目的性。为了\zd{公共利益}。
                            \item 适当性。\zd{手段有助于实现目的}。\al{为了控制狂犬病，把狗都杀了}。
                            \item 必要性。对民众权益\zd{侵害最小}。能不罚就不罚，能少罚就少罚。
                          \end{enumerate}
                  \end{enumerate}
            \item 程序正当。\zx{行政公开、公众参与、回避}。\al{注意听取\xhx{公民、法人 \hz 其他组织}的意见}。
            \item 高效\zx{便民}。便民不是便一两个民呐。\al{听取\xhx{群众/公众}意见}。
            \item 诚实守信。\zx{两信、两保}。
                  \begin{enumerate}
                    \item \zx{两信}。\yi 信息真实。别撒谎，别反悔。\er \zd{信赖利益保护}。防止朝令夕改、出尔反尔。
                    \item \zx{两保}。
                          \begin{enumerate}
                            \item \zd{存续保护}。非经法定事由，并由法定程序，不得撤销、变更已经生效了的行政决定。
                            \item 没办法存续 \yjt \zd{财产保护}。\al{赔偿 \cuo、补偿 \dui}。
                          \end{enumerate}
                  \end{enumerate}
            \item 权责统一。\zx{执法有保障、有权必有责、用权受监督、违法受追究、侵权需赔偿}。
                  \al{行政机关依法履行职责，法院、法规应赋予其相应的执法手段 \dui}。
          \end{enumerate}
          \tcblower
          \zx{多选题不能“多”选，即便表述对但不符合那个原则}。
        \end{kk}
\end{enumerate}
